\section{Detector and Evaluation Setup}\label{sec:setup}

\emph{CIC-IoT2023.} We use ten captures~\cite{neto2023ciciot2023}: Benign, DDoS-SYN, DoS-HTTP, Recon-PortScan, MITM-ArpSpoofing, DictionaryBruteForce, SqlInjection, XSS, CommandInjection and Mirai-udpplain (about 5.17\,GB). Files over 300\,MB are trimmed to their first $10^6$ packets. The captures map to eight classes: Benign, BruteForce, DDoS-SYN, DoS-HTTP, MITM-ARP (man-in-the-middle by ARP spoofing), Mirai, Recon and Web (SQL injection, XSS and command injection).

\emph{IoT-23.} We use ten IoT-23 captures~\cite{garcia2020iot23}: seven malware captures (Hide\&Seek, Muhstik, Hakai, two Torii, Mirai, Trojan) and three captures of real, uninfected IoT devices (Philips Hue bridge, Amazon Echo, Somfy smart door lock; of 7-1 only its Somfy-01 part), which we call benign-device captures. The Okiru captures (all at least 1\,GB) are left out. Labels are the dataset's own Zeek \texttt{conn.log.labeled}: each flow is joined to the Zeek connection with the same unordered 5-tuple and the nearest start within $\pm2$\,s, otherwise to the Zeek connection whose span contains the flow's start. Of 756,306 flows, 99.99\% match; the 38 unmatched flows are dropped, and only 270 (0.04\%) overlap Zeek rows of more than one class. The six classes are Benign, PortScan, C\&C, Attack (only in capture 3-1), DDoS and FileDownload.

\emph{Flows and features.} We build flows with NFStream~\cite{aouini2022nfstream} (version 6.6.0), using its statistical features, an idle timeout of 15\,s and an active timeout of 120\,s, with identical settings for both datasets and on the live gateway. The model uses 54 statistical features (packet counts and sizes, inter-arrival times, durations, TCP flag counts per direction). Identifiers (addresses, ports, protocol, timestamps) are excluded, because they let a model memorize the capture~\cite{arp2022dos}. CIC-IoT2023 yields 857,206 labeled flows, 511,527 of them DDoS-SYN (a consequence of trimming).

\emph{Detector.} The detector is a scikit-learn~\cite{pedregosa2011sklearn} decision tree with balanced class weights. Limiting depth and leaf size on validation raised the validation Benign FPR (from 1.3\% to 5.7\% at depth 16, leaf 10), so the \emph{deployed tree} (depth 32, leaf 1; 1,767 nodes, 222\,KiB) is effectively unconstrained. On the CIC-IoT2023 temporal split it scores macro-F1 0.725 (Benign FPR 3.8\%; 0.729 and 1.3\% with the label corrections of Section~\ref{sec:labels}), close to a random forest (0.737, at about 94$\times$ the size and 63$\times$ the single-flow latency) and histogram gradient boosting (0.714, Benign FPR 7.6\%). NSGA-II feature selection~\cite{deb2002nsga2,blank2020pymoo} kept 14 of 54 features at no measurable cost (0.723) but barely shrank the tree (211\,KiB). All results use the deployed tree, retrained with the same settings for each dataset and label set (except the check of the released CSVs in Section~\ref{sec:labels}).

\emph{Splits.} The \emph{temporal} split puts the first 70\% of each capture's flows, by start time, into training (on CIC-IoT2023, 600,050 training and 257,156 test flows). The \emph{random} split draws 70\% at random, stratified by class. Every design choice (tree limits, feature subset, blocking policy) is made on a validation set, the last 20\% of each capture's training period.

\emph{Metrics.} We report per-class F1 and recall with test support, macro-F1 over all classes, and the Benign false-positive rate (FPR), the share of Benign test flows predicted as any attack. A \emph{host pair} is an unordered pair of addresses within one capture. Flows from one host pair are not independent, so 95\% confidence intervals (CIs) come from a \emph{host-pair bootstrap}: 1,000 resamples of whole host pairs. A flow-level bootstrap~\cite{efron1994bootstrap} gives far too narrow intervals: [0.702, 0.741] for the temporal macro-F1, against [0.568, 0.822] with host pairs. Two splits are compared by a \emph{paired} bootstrap over the union of both test sets' host pairs: each resample draws host pairs with replacement and applies the same draw to both splits (a pair absent from one test set adds nothing to it).
